\documentclass[11pt]{article}
\usepackage[T1]{fontenc}
\usepackage[utf8]{inputenc}
\usepackage{lmodern,amsmath,amssymb,amsthm,geometry,hyperref}
\geometry{margin=1in}
\hypersetup{colorlinks=true,linkcolor=blue,urlcolor=blue,citecolor=blue}
\newtheorem{theorem}{Theorem}
\newtheorem{lemma}[theorem]{Lemma}
\theoremstyle{remark}\newtheorem{remark}[theorem]{Remark}
\newcommand{\ddc}{D}
\newcommand{\PSH}{\operatorname{PSH}}
\newcommand{\Eenv}{\mathcal E_{\mathrm{env}}}
\newcommand{\OO}{\mathcal O_X}
\title{Countable dense poles and equality of all obstacle envelopes}
\author{Research manuscript for Problem S09}
\date{9 October 2026\quad Private manuscript; explicit family independently checked}
\begin{document}
\maketitle
\begin{abstract}
For every positive Hermitian background and every prescribed countable
set, we construct genuine potentials with poles at all its points, trivial
multiplier ideals at every positive scale, and ordinary envelopes equal
to their ideal envelopes for every continuous obstacle. A fixed stronger
potential supplies ordinary-order competitors at every small convex
weight. Countable Jensen estimates prove exponential integrability
directly. A dense prescribed set gives potentials unbounded on every
nonempty open set. No mass bound or positive lower volume is assumed.
The conclusion concerns this explicit family and leaves the unrestricted
implication in original S09 unresolved. The explicit-family theorem has
passed independent reconstruction and a final scoped audit. Literature
priority is not asserted.
\end{abstract}

\section{Definitions and statement}
Let $X$ be connected compact complex, $\dim_{\mathbb C}X=n\geq1$, and
$h>0$ smooth Hermitian. Fix a smooth positive probability volume $dV$,
and use $\ddc=i\partial\bar\partial/(2\pi)$. All potentials have genuine
canonical quasi-psh representatives. Write $v\preceq u$ if $v\leq u+C$
for some finite constant. The original multiplier ideal $\mathcal I(su)$
uses the integrability of $|g|^2e^{-su}$, and $v\preceq_{\mathcal I}u$
means $\mathcal I(sv)\subseteq\mathcal I(su)$ for every real $s>0$.
For continuous real $F$ put
\[
 P_{h,F}[u]=\bigl(\sup\{v\in\PSH(X,h):v\leq F,
                                     \ v\preceq u\}\bigr)^*,
\]
with $P_{h,F}[u]_{\mathcal I}$ defined by ideal order instead, and
\[
 P_h(F)=\bigl(\sup\{v\in\PSH(X,h):v\leq F\}\bigr)^*.
\]
Here $*$ is upper semicontinuous regularization. These are the original
orders and envelopes, with no change of multiplier-ideal convention.

\begin{theorem}\label{th:main}
Fix $0<\beta<\beta'<1$ and any nonempty finite or countably infinite
set $\{x_j\}$ in $X$. There are genuine $u,w\in\PSH(X,h)$ satisfying:
\begin{enumerate}
\item $u,w\leq0$, $h+\ddc u,h+\ddc w\geq3h/4$, and
      $u(x_j)=w(x_j)=-\infty$ for every prescribed point.
\item For every $0<t\leq1$, $tw\preceq u$. The same fixed $w$
      works for all $t$. Also $u\in\Eenv(X,h,0)$ in the original
      stretching sense: for every $b\geq1$, some genuine $h$-psh
      $v$ satisfies $b^{-1}v\leq u$ canonically.
\item $e^{-su},e^{-sw}\in L^1(X,dV)$ for every real $s>0$;
      in particular $\mathcal I(su)=\OO$ for every such $s$.
\item For every continuous real obstacle $F$,
      \[
       P_{h,F}[u]=P_{h,F}[u]_{\mathcal I}=P_h(F).
      \]
      Consequently the two zero-obstacle envelopes both equal the
      genuine potential zero. Every original smooth obstacle is included.
\end{enumerate}
If the prescribed set is dense, $u$ is unbounded on every nonempty open
set. The pole sets can contain additional points. No Monge--Amp\`ere
mass or product-continuity assertion is made.
\end{theorem}

\section{Actual countable potentials}
Take a finite atlas of relatively compact coordinate charts covering $X$,
with slightly larger charts on which the coefficients of $dV$ are uniformly
bounded. For each $x_j$, choose a coordinate ball $B_j$ centered there
with radius $0<r_j\leq1$, contained in one of the larger charts. Let
$0\leq\chi\leq1$ be one smooth cutoff supported strictly inside the
unit ball and equal to one near zero. In that chart put
\[
 \chi_j(z)=\chi((z-z_j)/r_j),\qquad
 L_j=1-\log\frac{|z-z_j|^2}{r_j^2},
\]
and extend the supported profiles by zero outside $B_j$. The balls
may overlap. Let $e_j=2^{-j-2}$, so $\sum_j e_j\leq1/4$.

For $0<\sigma<1$, the inner profile $q_\sigma(s)=-(1-s)^\sigma$
has $q_\sigma'=\sigma(1-s)^{\sigma-1}>0$ and
$q_\sigma''=\sigma(1-\sigma)(1-s)^{\sigma-2}>0$.
Its tangential and radial Hessian eigenvalues in the normalized radial
coordinate are $q_\sigma'/R^2$ and $q_\sigma''/R^2$. The cutoff
transition is smooth with a finite negative Hessian bound relative to $h$.
Choose $\delta_j>0$, the same for $\beta$ and $\beta'$, with
$\delta_j\leq2^{-j}$ so that
\begin{equation}\label{eq:terms}
 a_j=-\delta_j\chi_j L_j^\beta,\qquad
 c_j=-\delta_j\chi_j L_j^{\beta'},\qquad
 \ddc a_j,\ddc c_j\geq-e_jh.
\end{equation}
This is possible separately for each $j$ by scaling both finite annular
bounds. At $x_j$ assign the canonical value $-\infty$.

For completeness this is genuine positivity across the point. Choose
smooth nondecreasing $\zeta$ equal to zero below $-2$ and to one above
$-1$. For $A$ large define
\[
 q_A'(s)=q_\sigma'(s)\zeta(s/A),\qquad
 q_A(s)=q_\sigma(0)-\int_s^0q_A'(t)\,dt,\quad s\leq0.
\]
The profile is convex increasing, constant for $s\leq-2A$, and equals
$q_\sigma$ for $s\geq-A$. It gives a smooth radial function at the
center and leaves the fixed cutoff transition unchanged. As $A$ increases
it decreases to $q_\sigma$, including the value $-\infty$ at the center,
and converges in $L^1$ by the radial integrability of $(1-s)^\sigma$.
The smooth Hessian bounds pass distributionally. Local addition of a
fixed strictly psh quadratic makes these a decreasing psh sequence, so
the limit has the claimed canonical representative
\cite[Theorems I.5.4 and I.5.8]{Dem}.

Define the actual decreasing sums
\[
 u=\lim_{N\to\infty}\sum_{j\leq N}a_j,\qquad
 w=\lim_{N\to\infty}\sum_{j\leq N}c_j.
\]
Finite-atlas volume bounds and rescaling by $r_j$ give
$\|a_j\|_{L^1}+\|c_j\|_{L^1}\leq C\delta_j$, since both unit-ball
radial integrals are finite. Thus the sums are $L^1$ and are not
identically $-\infty$. Each partial sum has background current at least
$(1-\sum_{j\leq N}e_j)h\geq3h/4$. A fixed local quadratic correcting
$h$ makes all partial sums psh. Their decreasing limits are therefore
genuine canonical quasi-psh functions, and their current bounds pass to
the limit. Each negative summand has a pole at its own $x_j$, proving
the asserted values. If these points are dense, every nonempty open set
contains an actual pole, so no local boundedness is possible there.

\section{One fixed stronger ordinary minorant}
Put $S=\sum_j\delta_j$, so $0<S<\infty$. For every real $b\geq1$,
\[
 M_b=\max\{0,\sup_{L\geq1}(bL^\beta-L^{\beta'})\}<\infty
\]
because $\beta'>\beta$. Termwise multiplication by $0\leq\chi_j\leq1$
and summation yield, for every finite $N$,
\[
 \sum_{j\leq N}c_j\leq b\sum_{j\leq N}a_j
                                  +M_b\sum_{j\leq N}\delta_j.
\]
Pass to the actual decreasing canonical sums. The inequality remains
valid at common poles; in particular if $u=-\infty$ it forces
$w=-\infty$. Hence $w\leq bu+M_bS$ everywhere. With $b=1/t$ this
gives
\begin{equation}\label{eq:minorant}
 tw\leq u+tM_{1/t}S,\qquad 0<t\leq1.
\end{equation}
The function $tw$ is genuine $h$-psh by convexity with the potential zero.
For stretching, $v=w-M_bS$ is genuine $h$-psh and satisfies $v\leq bu$.
Together with $u\leq0$, this proves the original $\Eenv(X,h,0)$ condition.
No bound uniform in $b$ or in $t$ on the ordinary-order constants is
required.

\section{Exponential integrability without valuation assumptions}
Let $s>0$ and $\alpha_j=\delta_j/S$. Weighted Jensen, assigning the
missing tail weight to the number zero, gives
\[
 e^{-s\sum_{j\leq N}a_j}
 \leq1-\sum_{j\leq N}\alpha_j
       +\sum_{j\leq N}\alpha_j e^{sS\chi_j L_j^\beta}.
\]
The individual integrals on the right are uniformly bounded in $j$.
Outside $B_j$ the exponential is one; inside, the finite-atlas volume
bound and $r_j^{2n}\leq1$ reduce the estimate to
\[
 C\int_0^1\exp\bigl(sS(1-\log R^2)^\beta\bigr)R^{2n-1}\,dR<\infty.
\]
With $x=-\log R^2$, the exponent is $-nx+sS(1+x)^\beta$,
which is eventually at most $-nx/2$ because $\beta<1$.
Integrating Jensen and then using monotone convergence proves
$e^{-su}\in L^1(X,dV)$. Repeat the same argument with $\beta'$ to
obtain $e^{-sw}\in L^1$ for every $s>0$. Every holomorphic germ is bounded
on a smaller coordinate ball. Thus $|g|^2e^{-su}$ is locally integrable
for every germ, proving $\mathcal I(su)=\OO$ in the exact original
convention. No selected-divisor or zero-Lelong-number inference is used.

\section{Every continuous obstacle}
\begin{lemma}\label{lem:envelope}
If a nonempty family of genuine $h$-psh functions is bounded above by
continuous $F$, its regularized supremum is genuine $h$-psh and below $F$.
\end{lemma}
\begin{proof}
First finite maxima retain the exact nonclosed bound. On a smaller chart,
convolve two potentials $v_i$ with smooth positive kernels. Positivity of
$h+\ddc v_i$ gives $\ddc(v_i*\kappa_\varepsilon)
\geq-h*\kappa_\varepsilon\geq-h-o(1)\ddc|z|^2$ uniformly there.
The smooth maximum
$m_\tau(x,y)=\tau\log(e^{x/\tau}+e^{y/\tau})$ has nonnegative first
derivatives summing to one and a positive semidefinite Hessian. Its
chain rule therefore preserves this lower bound. Convolution converges
in $L^1$, and $0\leq m_\tau-\max\leq\tau\log2$, so taking both limits
passes the exact bound $\ddc\max(v_1,v_2)\geq-h$ distributionally.
Adding one local quadratic $Q$ with $\ddc Q\geq h$ identifies its
canonical representative as $\max(v_1+Q,v_2+Q)-Q$
\cite[Theorem I.5.6]{Dem}.

Apply the local psh upper-envelope theorem and its Choquet selection
\cite[Theorems I.4.23 and I.5.7 and their proofs]{Dem} after adding the same $Q$.
A finite cover gives a countable subfamily whose supremum has the same
regularization and equals it almost everywhere. Include one fixed
genuine competitor $A$ in every finite maximum. These maxima increase,
lie between $A$ and $F$, and converge in $L^1$ by domination to the
regularized supremum. Their exact $h$ current bounds pass to that limit.
The local psh regularization supplies its canonical representative.
Continuity of $F$ preserves the upper bound at every point. No local
potential for the nonclosed form $h$ is required.
\end{proof}

Fix continuous real $F$ and set $W=P_h(F)$. The constant $\min_XF$
is a competitor, so $\min_XF\leq W\leq\max_XF$. Local psh supremum
regularization with Lemma~\ref{lem:envelope} gives a genuine bounded
$h$-psh envelope and preserves $W\leq F$.
Because every $\mathcal I(su)$ is $\OO$, every genuine $h$-psh function
is ideal-order dominated by $u$. Therefore $P_{h,F}[u]_{\mathcal I}=W$
directly from the full original definition.

For $0<t<1$ define the genuine convex-combination potential
\[
 A_t=(1-t)W+tw-t\|F\|_\infty.
\]
It is below $F$, since $w\leq0$, $W\leq F$, and
$(1-t)F-t\|F\|_\infty\leq F$. Equation \eqref{eq:minorant} and
boundedness of $W$ show $A_t\preceq u$ for each $t$. Hence $A_t$ is
an ordinary-order competitor for the original envelope $O=P_{h,F}[u]$.
The envelope is genuine and nontrivial because it dominates $A_t$;
all its competitors are bounded above by $F$, so
Lemma~\ref{lem:envelope} applies. We have $A_t\leq O\leq W$.
On the full-volume set where $w$ is finite, $A_t\to W$ as $t\downarrow0$,
so $O=W$ almost everywhere. Canonical equality holds everywhere: locally
add one smooth quadratic to both, obtaining two psh functions with the
same distribution; uniqueness of the psh representative
\cite[Theorem I.5.8]{Dem} makes them identical, including point values.
This proves the asserted equality for all $F$. For $F=0$, $W=0$ because
zero itself is a competitor.

\begin{remark}
These statements are independent of any bounded upper mass property or
positive lower singularity volume. They apply on any given positive
background having zero lower mass or unbounded upper mass; no such
geometry or top-product theorem is inferred from the construction.
Original S09 concerns every genuine zero-obstacle equality for arbitrary
admissible smooth backgrounds. Here positivity of $h$ and the fixed
stronger minorant are explicit restrictions. The known relative-envelope
sufficient criterion is not claimed as a new universal theorem. The
countable construction and all of its elementary primary applications
have passed independent reconstruction and final scoped audit. This
acceptance covers only the explicit family stated here.
\end{remark}

\begin{thebibliography}{9}
\bibitem{Dem} J.-P. Demailly, \emph{Complex Analytic and Differential
Geometry}, author-hosted manuscript, version of 21 June 2012, Chapter I,
Theorems 4.23, 5.4, 5.6, 5.7 and 5.8.
\end{thebibliography}
\end{document}
